\section*{Chapitre \ref{ch:s2:tail-hedging-and-long-volatility} --- Couverture de queue et volatilité longue}

\begin{solution}{exo:s2:tail-hedging-and-long-volatility:1}
$17 \times (-6.0) + 19.2 - 2.7 + 39.5 = -46.0~\%$ du portefeuille, additionnés sans capitalisation~: la couverture a perdu de l'argent au cours des vingt années, et la
croissance du portefeuille est tombée de 3.0~\% à 2.5~\% par an, tandis que son pire drawdown est passé de 73~\% à 41~\%.
\end{solution}

\begin{solution}{exo:s2:tail-hedging-and-long-volatility:2}
Le put de chaque mois est fixé à 5~\% sous le niveau de départ de ce mois, et chaque mois l'indice baisse d'environ 4~\%~: aucun put ne finit dans la monnaie, et chaque
prime est perdue, alors que l'indice baisse d'environ 40~\% sur l'année.
\end{solution}

\begin{solution}{exo:s2:tail-hedging-and-long-volatility:3}
Trois fois son coût~: 3~\% du notionnel.
\end{solution}

\begin{solution}{exo:s2:tail-hedging-and-long-volatility:4}
Le put spread de CLLZ ne protège qu'entre 95~\% et 97.5~\% du niveau~; une baisse de 30~\% traverse le put spread, et le call vendu abandonne tout
rebond. Le put unique de PPUT protège sur toute la baisse.
\end{solution}

\begin{solution}{exo:s2:tail-hedging-and-long-volatility:5}
PPUT est un indice de rendement total qui réinvestit les dividendes~; l'indice de cours S\&P 500 ne le fait pas. Le véritable indice de référence croît plus vite que l'indice de cours, donc
le manque à gagner réel de l'indice couvert est supérieur à l'écart affiché.
\end{solution}

\begin{solution}{exo:s2:tail-hedging-and-long-volatility:6}
Une surcouche de tendance ne se met en position courte qu'après que le rendement glissant est devenu négatif~: elle protège lors des longues baisses et ne répond à aucun krach de quelques jours. Un
put protège dès son prix d'exercice, quel que soit le chemin, et coûte une prime, qu'un krach survienne ou non.
\end{solution}

\begin{solution}{exo:s2:tail-hedging-and-long-volatility:7}
Sur cette année, l'indice perd 28.7~\%, l'allocation statique 21.0~\% et le programme de couverture par puts 32.4~\%, davantage que l'indice non couvert. Le put de chaque mois est fixé
à 5~\% sous le début du mois, et l'indice baisse de moins que cela au cours de la plupart des mois, si bien que les puts rapportent rarement et que leurs primes, plus chères à mesure que la volatilité monte,
aggravent la perte. C'est le point d'Israelov dans le modèle du chapitre.
\end{solution}

\begin{solution}{exo:s2:tail-hedging-and-long-volatility:8}
Cinq années calmes ne disent rien d'une couverture conçue pour des krachs rares~: sur le marché synthétique, le programme a perdu de l'argent pendant 17 années sur 20 et a rapporté 19.2~\% et
39.5~\% dans deux des autres années. Il faut la juger à l'aune du portefeuille, sur un historique comportant des krachs, et décider à l'avance si son coût achète une exposition supplémentaire
ailleurs.
\end{solution}

\begin{solution}{pb:s2:tail-hedging-and-long-volatility:1}
\begin{enumerate}
\item Une position qui paie lors de fortes baisses rares, achetée à un coût~; son \omterm{def:s2:tail-hedging-and-long-volatility:bleed}{coût de portage de la couverture} (bleed) est ce coût pendant les périodes calmes, rapporté au portefeuille par an.
\item La prime de variance sur la volatilité à la monnaie et la prime de skew des puts hors de la monnaie.
\item Les puts à courte échéance sont la défense la plus fiable et la plus coûteuse~; les obligations sont peu fiables~; la tendance et la qualité ont bien résisté lors des drawdowns.
\item Les puts ne valent souvent pas mieux que la détention statique d'un portefeuille moins exposé, sauf si leur moment d'achat est bien choisi.
\item PPUT 7.7~\% pour une volatilité de 13.5~\%, CLLZ 7.8~\% pour 14.6~\%, l'indice de cours S\&P 500 8.5~\% pour 18.4~\% (sans dividendes).
\item 1987~: $-10.6~\%$, $-17.1~\%$, $-21.8~\%$~; automne 2008~: $-6.7~\%$, $-31.2~\%$, $-30.1~\%$~; février--mars 2020~: $-0.4~\%$, $-17.9~\%$, $-19.9~\%$.
\item La baisse a duré vingt mois~; le put de chaque mois n'a rapporté que pour la baisse de ce mois au-delà de 5~\%.
\item Il ne protège qu'entre ses prix d'exercice.
\item Vendre une couverture qui a gagné avant l'échéance et la remplacer au nouveau niveau.
\item D'après le tableau~: puts mensuels à 5~\%, 2.5~\% par an avec un drawdown de $-41~\%$~; puts semestriels à 20~\%, 4.6~\% et $-50~\%$~; couverture monétisée, 2.2~\% et $-62~\%$~;
l'indice 3.0~\% et $-73~\%$.
\item Un coût de portage de 6.0~\% par an pendant 17 années calmes~; 19.2~\%, $-2.7~\%$ et 39.5~\% lors des années de krach.
\item Le put de remplacement, fixé près du nouveau niveau inférieur, a laissé le reste de la baisse sans protection.
\item \textbf{Résultat à retenir.} Le programme mensuel à 5~\% coûte 6.0~\% du portefeuille par an pendant les années calmes~; le portefeuille couvert croît de 2.5~\% par an
contre 3.0~\% sans couverture, le pire drawdown passant de 73~\% à 41~\%.
\item Volatilité voisine (13.9~\% contre 13.2~\%), drawdown plus faible ($-41~\%$ contre $-58~\%$), croissance plus faible (2.5~\% contre 2.8~\%).
\item La surcouche de tendance croît plus vite (6.3~\%), avec un drawdown plus profond ($-56~\%$).
\item Ses krachs durent dix jours, le cas pour lequel les puts sont conçus~; les marchés baissiers lents sont absents.
\item Par le taux de croissance du portefeuille sur les krachs des deux types, et par l'exposition supplémentaire que permet sa protection.
\item Les drawdowns rapides et lents, les prix des options sur le smile, une alternative statique, la monétisation aux prix de l'époque.
\item La tendance (ou détenir moins).
\item La convexité lors d'un krach, payée dans tous les autres mois.
\end{enumerate}
\end{solution}

\begin{solution}{iq:s2:tail-hedging-and-long-volatility:1}
Les acheteurs veulent une protection contre les krachs et les vendeurs doivent être rémunérés pour la fournir~: les puts sur indice portent la prime de variance et la prime de skew, toutes deux maximales pour
les puts sur indice hors de la monnaie.
\end{solution}

\begin{solution}{iq:s2:tail-hedging-and-long-volatility:2}
Égaliser les deux stratégies sur la volatilité ou sur le drawdown attendu, puis comparer les taux de croissance et les drawdowns sur un historique comportant des baisses rapides et lentes, comme l'a fait Israelov~;
inclure les coûts et le calendrier des roll.
\end{solution}

\begin{solution}{iq:s2:tail-hedging-and-long-volatility:3}
Suivre la règle fixée à l'avance. Vendre fige le gain en vue d'un réinvestissement et réinitialise la protection près du nouveau niveau~; conserver les puts maintient la
protection contre une nouvelle baisse. Un compromis courant consiste à vendre une partie et à reporter le reste vers des prix d'exercice plus bas.
\end{solution}

\begin{solution}{iq:s2:tail-hedging-and-long-volatility:4}
Un krach rapide, une baisse lente sur un an, un krach juste après un roll et un autre juste avant, un pic de volatilité sans baisse, et plusieurs années calmes
d'affilée pour tester le budget.
\end{solution}

\begin{solution}{iq:s2:tail-hedging-and-long-volatility:5}
Chaque jour~: valoriser les puts sur la surface de volatilité, vérifier la règle de monétisation par rapport au coût enregistré à l'achat, faire les roll selon le calendrier, dimensionner les nouveaux puts en fonction du
portefeuille, enregistrer le \omterm{def:s2:tail-hedging-and-long-volatility:bleed}{coût de portage de la couverture}~; avec des contrôles sur les cotations d'options périmées et sur les dates de roll.
\end{solution}

\begin{solution}{iq:s2:tail-hedging-and-long-volatility:6}
La richesse en $\log$ croît de $\ln(1 + R)$ à chaque période et $R$ est log-normal, de moyenne $\mu$ et de variance $\sigma^2$, $E[\ln(1 + R)] \approx \mu -
\sigma^2/2$. Une couverture qui abaisse $\mu$ de $c$ et $\sigma^2$ de plus de $2c$ relève le taux de croissance~: c'est le cas d'une couverture qui permet aussi une
position plus grande sur l'actif risqué.
\end{solution}
